% Related Work Section for STAR Whitepaper
% Insert after Section 2 (Mathematical Foundation) and before Section 3 (System Architecture)

\section{Related Work}
\label{sec:related}

\subsection{Vector-Based Retrieval-Augmented Generation}

Modern RAG systems predominantly rely on dense vector representations and approximate nearest neighbor (ANN) search. HNSW (Hierarchical Navigable Small World) graphs \cite{malkov2018efficient} and FAISS \cite{johnson2019billion} represent the state-of-the-art for vector retrieval, offering sub-linear query complexity. However, these approaches require loading complete indices into RAM-often gigabytes for modest corpora-restricting deployment to high-specification servers. Furthermore, vector similarity provides limited explainability: a result matches because its embedding is "close" to the query, but the specific reasoning remains opaque. STAR addresses these limitations through sparse graph traversal, enabling CPU-only deployment on resource-constrained devices while providing explicit tag-based provenance for every result.

\subsection{Graph-Based Memory Systems}

Recent work has explored graph structures as alternatives to dense vectors. T-Retriever \cite{wei2026tretriever} introduces tree-based hierarchical retrieval using semantic-structural entropy for encoding textual graphs. While effective for hierarchical document structures, T-Retriever does not incorporate temporal decay-a key requirement for personal memory systems where recency matters. PersonalAI \cite{menschikov2025personalai} proposes a knowledge graph framework with hyper-edges for personalized LLM agents, achieving strong results on TriviaQA and HotpotQA benchmarks. However, PersonalAI focuses on framework design rather than production implementation; STAR contributes a complete, deployed system with validated performance on 28M tokens of real-world data.

Our bipartite graph approach (Atoms $\times$ Tags) differs from general knowledge graphs by enforcing a strict separation between content and metadata. This enables O(1) deduplication via SimHash \cite{charikar2002similar} and supports disposable index architectures where the database can be rebuilt entirely from the source-of-truth filesystem.

\subsection{Personal AI Memory Systems}

The advent of large context windows has renewed interest in personal AI memory. Second Me \cite{wei2025second} proposes LLM-based memory parameterization, using language models themselves to structure and retrieve personal knowledge. While powerful, this approach requires significant computational resources and offers limited explainability. STAR achieves similar associative retrieval goals through deterministic physics-based scoring, enabling deployment on 4GB RAM laptops without GPU acceleration.

Cognitive AI frameworks \cite{salas2025cognitive} emphasize governed memory architectures for long-term coherence. STAR's ephemeral index design (Standard 110) aligns with these principles while adding practical constraints for local-first deployment: zero cloud dependencies, AGPL-3.0 licensing, and real-world validation.

\subsection{Temporal Information Retrieval}

Temporal decay has been explored in web archive search \cite{kanhabua2008surviving} and recency-weighted ranking, but is rarely integrated into RAG systems as a fundamental scoring component. STAR's Unified Field Equation (Equation~\ref{eq:unified_field}) embeds temporal decay multiplicatively alongside semantic and structural factors, ensuring that any zero factor eliminates irrelevant results. This differs from additive scoring approaches where weak signals can accumulate noise.

\subsection{Local-First and Edge Computing}

The local-first software movement \cite{haque2023local} emphasizes user data ownership and offline capability. STAR's browser paradigm extends these principles to AI memory: just as browsers render content without downloading the entire internet, STAR retrieves context without loading complete vector indices. This enables sovereign operation-users maintain complete control over their data without cloud dependencies.

\subsection{Summary of Contributions}

STAR distinguishes itself from prior work through:
\begin{enumerate}
    \item \textbf{Sparse Graph Physics:} Multiplicative scoring combining co-occurrence, temporal decay, and SimHash similarity (Section \ref{sec:math}).
    \item \textbf{Browser Paradigm:} Sharded atomization enabling resource-constrained devices to navigate large corpora (Section \ref{sec:architecture}).
    \item \textbf{Production Validation:} Real-world deployment with 28M tokens, $<$200ms p95 latency, and 4GB RAM compatibility (Section \ref{sec:benchmarks}).
    \item \textbf{Explainable Retrieval:} Tag paths provide deterministic provenance for every result (Section \ref{sec:retrieval}).
\end{enumerate}